

\hypertarget{fixture-chicago17-nb-parity}{%
\section{Fixture: chicago17-nb
parity}\label{fixture-chicago17-nb-parity}}

Fernández opens the discussion with a sweeping claim about liturgical
continuity across the early medieval West.\footnote{María Luisa
  Fernández, \emph{Liturgical Continuity in the Early Medieval West}
  (Cambridge: Cambridge University Press, 2011).} That framing set the
terms for a generation of later work, including the close palaeographic
study produced in Leipzig a decade later.\footnote{Klaus Herrmann and
  Annika Schulz, {``A Palaeographic Reassessment of the Leipzig
  Codices,''} \emph{Speculum} 96, no. 2 (2021): 38--71.}

Fernández returns to the same argument in a later chapter, now with a
harder quantitative edge:\footnote{Fernández, \emph{Liturgical
  Continuity in the Early Medieval West}, 214.}

\begin{quote}
The codices in question share not only a scribal hand but a deliberate
editorial program, and any account that treats them as independent
witnesses misreads the material evidence on its face.\footnote{Fernández,
  219.}
\end{quote}

Collaborative work on the same corpus broadened the claim to four
regional scriptoria,\footnote{Chidinma Okafor et al., {``Four
  Scriptoria, One Editorial Program: A Comparative Study,''}
  \emph{Scriptorium} 72, no. 1 (2018): 5--48.} and Herrmann and Schulz
returned to the question in a short follow-up note.\footnote{Herrmann
  and Schulz, {``A Palaeographic Reassessment of the Leipzig Codices,''}
  44.} A more recent archival guide fills in the institutional
background.\footnote{Bibliothèque nationale de France, \emph{Guide to
  the Medieval Manuscript Collections} (Paris: Bibliothèque nationale de
  France, 2022), \url{https://www.bnf.fr/fr/guide-manuscrits-medievaux}.}

\hypertarget{bibliography}{%
\section*{Bibliography}\label{bibliography}}
\addcontentsline{toc}{section}{Bibliography}

\hypertarget{refs}{}
\begin{CSLReferences}{1}{0}
\leavevmode\vadjust pre{\hypertarget{ref-bnf-guide-2022}{}}%
Bibliothèque nationale de France. \emph{Guide to the Medieval Manuscript
Collections}. Paris: Bibliothèque nationale de France, 2022.
\url{https://www.bnf.fr/fr/guide-manuscrits-medievaux}.

\leavevmode\vadjust pre{\hypertarget{ref-fernandez-2011}{}}%
Fernández, María Luisa. \emph{Liturgical Continuity in the Early
Medieval West}. Cambridge: Cambridge University Press, 2011.

\leavevmode\vadjust pre{\hypertarget{ref-herrmann-schulz-2021}{}}%
Herrmann, Klaus, and Annika Schulz. {``A Palaeographic Reassessment of
the Leipzig Codices.''} \emph{Speculum} 96, no. 2 (2021): 38--71.

\leavevmode\vadjust pre{\hypertarget{ref-okafor-lindqvist-tanaka-rossi-2018}{}}%
Okafor, Chidinma, Erik Lindqvist, Hiroshi Tanaka, and Giulia Rossi.
{``Four Scriptoria, One Editorial Program: A Comparative Study.''}
\emph{Scriptorium} 72, no. 1 (2018): 5--48.

\end{CSLReferences}

